\chapter{Resultados}\label{cap_4}

Este capítulo apresenta os resultados em uma ordem que acompanha o raciocínio da investigação, e não a cronologia dos experimentos. Primeiro a reprodução do trabalho de referência, que define o ponto de partida. Depois a avaliação dos critérios de \ac{AL} sobre a seleção de imagens, que é a forma usual de aplicar a técnica. Em seguida o diagnóstico que mede quanta margem existe na escolha da região, e a caracterização de para onde cada critério aponta, que juntos explicam por que os critérios conhecidos não alcançam essa margem. Os componentes de \ac{CL} e \ac{IL} vêm na sequência. O capítulo fecha com a intervenção derivada do diagnóstico e com os resultados de natureza metodológica.

Todos os números vêm do registro de execuções descrito no Capítulo \ref{cap_3}. Cada tabela é gerada a partir desse registro por um programa, sem passo manual, e acompanha um arquivo de procedência que lista os identificadores das execuções que produziram cada linha. Nenhum valor foi transcrito à mão.

\section{Reprodução do trabalho de referência}

A reprodução do arranjo descrito por \cite{soares2026flim} confirma o comportamento central relatado pelos autores: o $F_\beta$ cresce de forma monotônica à medida que imagens são acrescentadas ao conjunto de treino e satura entre três e quatro imagens, atingindo valores próximos aos publicados para o codificador.

Duas diferenças de implementação precisam ser declaradas antes de qualquer comparação numérica. A primeira é que a avaliação usa limiarização de Otsu seguida de filtro de área, e não o delineamento por árvores dinâmicas empregado no trabalho original, cujo executável é um binário compilado para Linux que não roda no ambiente usado aqui. A segunda é que a seleção de imagens do Algoritmo 1 do trabalho de referência exige o ground truth de todo o conjunto disponível para escolher as poucas imagens que serão anotadas, o que os próprios autores declaram ao descrever a abordagem como supervisionada.

Essa segunda diferença não é um detalhe de implementação, e sim o ponto de partida desta dissertação. Se a escolha das imagens depende de conhecer a resposta em todas elas, o procedimento não é aplicável na situação que motiva o \ac{FLIM}, que é justamente a ausência de anotação. A pergunta que organiza os experimentos seguintes é quanto do benefício dessa escolha se recupera com critérios que não usam ground truth algum.

A Figura~\ref{fig_curva} e a Tabela~\ref{tab_curva} mostram o comportamento da curva de orçamento. A saturação entre três e quatro imagens é o que responde à pergunta de por que o trabalho de referência usa esse número: o algoritmo para quando acrescentar imagem deixa de melhorar o $F_\beta$.

\begin{figure}[htb]
\centering
\includegraphics[width=0.86\linewidth]{curva_orcamento.png}
\caption[Qualidade em função do número de imagens anotadas]{Qualidade em função do número de imagens anotadas, com a faixa de ruído da reexecução. A curva satura entre três e quatro imagens, e o ganho posterior cabe dentro do intervalo de confiança.}
\label{fig_curva}
\end{figure}

\begin{table}[htb]
\centering
\caption{Qualidade por número de imagens anotadas.}
\label{tab_curva}
% Fβ por número de imagens anotadas
% GERADO por flim_al/tabelas.py em 2026-10-08T14:01:05+00:00
% NAO EDITE A MAO: a proxima geracao sobrescreve, e o numero editado deixa de bater com evidencia/execucoes/.
% Proveniencia: curva_orcamento.proveniencia.json
\begin{tabular}{lrrrrr}
\toprule
imagens & F$\beta$ médio & desvio & IC95\% & n & seeds \\
\midrule
1 & 0.602 & 0.084 & [0.580, 0.625] & 54 & 9 \\
2 & 0.576 & 0.143 & [0.538, 0.615] & 54 & 9 \\
3 & 0.645 & 0.074 & [0.625, 0.664] & 57 & 9 \\
4 & 0.670 & 0.077 & [0.650, 0.691] & 54 & 9 \\
5 & 0.665 & 0.093 & [0.642, 0.688] & 63 & 9 \\
6 & 0.645 & 0.104 & [0.619, 0.672] & 58 & 9 \\
7 & 0.680 & 0.080 & [0.659, 0.701] & 54 & 9 \\
8 & 0.676 & 0.085 & [0.656, 0.697] & 65 & 9 \\
9 & 0.686 & 0.063 & [0.671, 0.702] & 64 & 9 \\
10 & 0.599 & 0.111 & [0.554, 0.644] & 26 & 5 \\
12 & 0.767 & 0.000 & [0.767, 0.767] & 3 & 0 \\
16 & 0.781 & 0.002 & [0.778, 0.785] & 3 & 0 \\
20 & 0.772 & 0.001 & [0.771, 0.773] & 3 & 0 \\
25 & 0.759 & 0.006 & [0.743, 0.775] & 3 & 0 \\
31 & 0.711 & 0.000 & [0.711, 0.711] & 3 & 0 \\
\bottomrule
\end{tabular}

\end{table}

% Nota de curva_orcamento. GERADA por flim_al/tabelas.py.
\footnotesize Seleção sem consulta ao gabarito. \texttt{n} conta execuções; \texttt{seeds} conta sementes distintas --- nove execuções de uma semente só não são nove repetições.
\normalsize


Uma ressalva importante acompanha essa leitura, e ela diz respeito ao que a curva mistura. A arquitetura pede duzentos núcleos por camada, mas o agrupamento só produz tantos quantos os trechos sob os marcadores permitem. Com uma imagem anotada o codificador sai com 54, 51, 48 e 48 núcleos nas quatro camadas; com oito, sai com duzentos em todas. Com pouca anotação a rede não é apenas menos treinada, \textbf{ela é menor}. Parte do ganho que a curva mostra nos primeiros pontos é, portanto, aumento de capacidade, e não aumento de supervisão, e separar as duas contribuições exigiria fixar o tamanho do codificador ao longo da curva.

Nas cinco últimas linhas da Tabela~\ref{tab_curva}, referentes a doze imagens ou mais, a coluna de sementes distintas indica zero. Isso significa que o registro não guardou semente distinta para essas execuções, de modo que três execuções de uma mesma semente não constituem três repetições, e o desvio nulo é artefato disso, não estabilidade.

A reimplementação do Algoritmo 1 revelou ainda duas propriedades não discutidas no trabalho original. O critério de selecionar a imagem de pior $F_\beta$ tende a escolher imagens insegmentáveis, e não imagens informativas: as imagens com $F_\beta$ igual a zero são aquelas cujas componentes previstas caem fora da faixa aceita pelo filtro de área, de modo que o algoritmo prioriza casos em que nenhuma anotação ajudaria. Além disso, sem uma lista de imagens rejeitadas, a imagem removida pelo retrocesso volta a ser a de menor $F_\beta$ na iteração seguinte, e o procedimento pode entrar em ciclo.

\section{Seleção de imagens}

A primeira pergunta é se algum critério de \ac{AL} sem ground truth supera a escolha aleatória das imagens a anotar. Foram avaliados critérios das duas famílias predominantes na literatura: incerteza, representada por entropia, \textit{least confidence}, margem e BALD; e diversidade, representada por CoreSet, BADGE e medoide. A referência superior é o oracle, que usa o ground truth para escolher a imagem de pior desempenho, e a referência inferior é o sorteio.

A Tabela \ref{tab_geral_al} resume a comparação nos três conjuntos de dados, para quatro orçamentos de anotação.

\begin{table}[htb]
\centering
\caption{Critérios de Active Learning por orçamento e conjunto de dados.}
\label{tab_geral_al}
% Active Learning no FLIM: todos os critérios, todos os K (FLIM_lm)
% GERADO por flim_al/tabelas.py em 2026-10-06T16:49:56+00:00
% NAO EDITE A MAO: a proxima geracao sobrescreve, e o numero editado deixa de bater com evidencia/execucoes/.
% Proveniencia: geral_al_lm.proveniencia.json
\begin{tabular}{llrrrrrr}
\hline
critério & K & Parasitas & BraTS & Conjuntivite & $\Delta$ médio vs sorteio & p & vence o sorteio em \\
\hline
CoreSet (imagem) & 2 & 0.543 & 0.000 & 0.642 & -0.061 & 0.5568 & 3 de 9 \\
CoreSet (imagem) & 3 & 0.538 & 0.033 & 0.587 & 0.023 & 0.6074 & 3 de 9 \\
CoreSet (imagem) & 5 & 0.640 & 0.489 & 0.748 & 0.112 & 0.3323 & 7 de 9 \\
CoreSet (imagem) & 8 & 0.690 & 0.230 & 0.636 & -0.067 & 0.4144 & 3 de 9 \\
Entropia (imagem) & 2 & 0.457 & 0.001 & 0.635 & -0.092 & 0.3677 & 3 de 9 \\
Entropia (imagem) & 3 & 0.474 & 0.000 & 0.591 & -0.008 & 0.9104 & 2 de 9 \\
Entropia (imagem) & 5 & 0.599 & 0.498 & 0.667 & 0.074 & 0.5274 & 6 de 9 \\
Entropia (imagem) & 8 & 0.681 & 0.000 & 0.604 & -0.157 & 0.0427 & 2 de 9 \\
Least confidence (imagem) & 2 & 0.420 & 0.001 & 0.635 & -0.105 & 0.3115 & 3 de 9 \\
Least confidence (imagem) & 3 & 0.451 & 0.000 & 0.591 & -0.016 & 0.8178 & 2 de 9 \\
Least confidence (imagem) & 5 & 0.583 & 0.498 & 0.667 & 0.069 & 0.5531 & 6 de 9 \\
Least confidence (imagem) & 8 & 0.684 & 0.000 & 0.441 & -0.210 & 0.0327 & 2 de 9 \\
AL de REGIAO & 2 & 0.142 & 0.252 & 0.442 & -0.178 & 0.2461 & 1 de 9 \\
AL de REGIAO & 3 & 0.280 & 0.257 & 0.443 & -0.036 & 0.7332 & 3 de 9 \\
AL de REGIAO & 5 & 0.039 & 0.499 & 0.261 & -0.247 & 0.1570 & 2 de 9 \\
AL de REGIAO & 8 & 0.412 & 0.525 & 0.308 & -0.170 & 0.1903 & 2 de 9 \\
Oracle (usa GT) & 2 & 0.093 & 0.196 & 0.274 & -0.269 & 0.0921 & 1 de 9 \\
Oracle (usa GT) & 3 & 0.420 & 0.630 & 0.312 & 0.091 & 0.5347 & 5 de 9 \\
Oracle (usa GT) & 5 & 0.618 & 0.445 & 0.572 & 0.031 & 0.7792 & 6 de 9 \\
Oracle (usa GT) & 8 & 0.518 & 0.008 & 0.535 & -0.232 & 0.0107 & 0 de 9 \\
\hline
\end{tabular}

\end{table}

% Nota de geral_al_lm. GERADA por flim_al/tabelas.py.
\footnotesize Uma linha por (critério, K). As colunas de dataset são o F$\beta$ médio sobre as sementes. \textbf{\texttt{FLIM puro (sorteio)} é a referência}: todos os braços recebem o MESMO número de imagens e são medidos no mesmo conjunto de teste. Os critérios de IMAGEM escolhem quais imagens anotar; \textbf{\texttt{AL de REGIÃO} usa as mesmas imagens do sorteio e muda só ONDE o traço cai}, com a mesma contagem de pixels --- o $\Delta$ dele mede posição de anotação, não escolha de imagem. \textbf{\texttt{Oracle} usa o gabarito} para escolher a imagem em que o modelo vai pior (passo 7 do Algoritmo 1 do artigo): não é método, é referência, e o fato de ele também não se destacar é o achado mais informativo daqui. $\Delta$ e p são pareados por semente dentro de cada dataset e K, depois agregados. \texttt{vence o sorteio em} conta as células (dataset $\times$ semente) favoráveis. Decoder FLIM\_lm.
\normalsize


Nenhum critério apresenta superioridade sobre o sorteio que sobreviva à correção de Benjamini-Hochberg. O resultado se mantém nos três conjuntos de dados e nos quatro orçamentos. Vale notar que o oracle, que lê o ground truth e portanto representa o teto do que a seleção de imagens poderia entregar, também não se destaca de forma consistente, perdendo do sorteio em parte considerável das células. Isso sugere que, no regime de três a cinco imagens, a identidade das imagens escolhidas importa menos do que a formulação do Algoritmo 1 pressupõe.

É importante enunciar esse resultado com precisão. O que foi observado é ausência de evidência de superioridade nas condições avaliadas, e não demonstração de equivalência. Os intervalos de confiança são compatíveis com efeitos pequenos em qualquer direção, e o número de repetições limita o tamanho de efeito detectável.

A Tabela~\ref{tab_comparacao_criterios} apresenta o conjunto completo de técnicas que chegaram a rodar, incluindo as que perderam do sorteio. Ela é deliberadamente exaustiva: apresentar apenas as técnicas que se saíram melhor seria seleção de resultado, e a leitura correta deste capítulo depende de ver que o fenômeno é geral e não um acidente de uma implementação.

\begin{table}[htb]
\centering
\caption{Todas as técnicas avaliadas, contra o sorteio.}
\label{tab_comparacao_criterios}
% Técnicas de Active Learning contra o sorteio
% GERADO por flim_al/tabelas.py em 2026-10-08T14:01:04+00:00
% NAO EDITE A MAO: a proxima geracao sobrescreve, e o numero editado deixa de bater com evidencia/execucoes/.
% Proveniencia: comparacao_criterios.proveniencia.json
\begin{adjustbox}{max width=\linewidth}
\begin{tabular}{lrrrrrrrrr}
\toprule
técnica & nível & F$\beta$ (células pareadas) & n pares & $\Delta$ vs random & p & $\Delta$ vs FLIM (K=3) & p (FLIM) & $\Delta$ vs controle do nível & veredito \\
\midrule
badge & imagem & 0.430 & 95 & \textbf{-0.019} & 0.0105 & --- & --- & --- & suportado \\
coreset & imagem & 0.464 & 354 & \textbf{-0.028} & 0.0149 & --- & --- & --- & suportado \\
entropy & imagem & 0.455 & 471 & -0.009 & 0.3661 & --- & --- & --- & sem suporte \\
least\_confidence & imagem & 0.437 & 369 & -0.020 & 0.0937 & --- & --- & --- & indicativo \\
margin & imagem & 0.451 & 72 & 0.023 & 0.2110 & --- & --- & --- & sem suporte \\
random & imagem & --- & --- & --- & --- & --- & --- & --- & piso \\
oracle & imagem (usa GT) & 0.514 & 103 & -0.020 & 0.4456 & --- & --- & --- & sem suporte \\
flim\_3img & linha de base & --- & 0 & --- & --- & --- & --- & --- & sem pares suficientes \\
random\_region & região & 0.442 & 252 & \textbf{0.108} & 0.0000 & --- & --- & --- & suportado \\
region\_bald & região & 0.462 & 135 & \textbf{0.077} & 0.0000 & --- & --- & 0.050 & suportado \\
region\_entropy & região & 0.437 & 66 & 0.014 & 0.4295 & --- & --- & --- & sem suporte \\
1\_imagens & ? & --- & 0 & --- & --- & --- & --- & --- & sem pares suficientes \\
2\_imagens & ? & --- & 0 & --- & --- & --- & --- & --- & sem pares suficientes \\
3\_imagens & ? & --- & 0 & --- & --- & --- & --- & --- & sem pares suficientes \\
5\_imagens & ? & --- & 0 & --- & --- & --- & --- & --- & sem pares suficientes \\
8\_imagens & ? & --- & 0 & --- & --- & --- & --- & --- & sem pares suficientes \\
al\_regiao & ? & 0.456 & 28 & \textbf{-0.032} & 0.0439 & --- & --- & --- & suportado \\
aleatorio\_objeto & ? & --- & 0 & --- & --- & --- & --- & --- & sem pares suficientes \\
argmax\_entropia & ? & --- & 0 & --- & --- & --- & --- & --- & sem pares suficientes \\
argmax\_lc & ? & --- & 0 & --- & --- & --- & --- & --- & sem pares suficientes \\
argmax\_prob & ? & --- & 0 & --- & --- & --- & --- & --- & sem pares suficientes \\
artigo & ? & --- & 0 & --- & --- & --- & --- & --- & sem pares suficientes \\
base & ? & --- & 0 & --- & --- & --- & --- & --- & sem pares suficientes \\
contorno\_sorteado & ? & --- & 0 & --- & --- & --- & --- & --- & sem pares suficientes \\
flim\_paper & ? & 0.467 & 69 & -0.015 & 0.4074 & --- & --- & --- & sem suporte \\
flim\_puro & ? & --- & 0 & --- & --- & --- & --- & --- & sem pares suficientes \\
il\_alpha0.00 & ? & --- & 0 & --- & --- & --- & --- & --- & sem pares suficientes \\
il\_alpha0.50 & ? & --- & 0 & --- & --- & --- & --- & --- & sem pares suficientes \\
il\_alpha1.00 & ? & --- & 0 & --- & --- & --- & --- & --- & sem pares suficientes \\
imagem\_nova & ? & --- & 0 & --- & --- & --- & --- & --- & sem pares suficientes \\
medoide & ? & 0.482 & 94 & \textbf{-0.043} & 0.0054 & --- & --- & --- & suportado \\
melhor\_sorteada & ? & --- & 0 & --- & --- & --- & --- & --- & sem pares suficientes \\
pior\_sorteada & ? & --- & 0 & --- & --- & --- & --- & --- & sem pares suficientes \\
regiao\_aleatoria & ? & --- & 0 & --- & --- & --- & --- & --- & sem pares suficientes \\
regiao\_borda & ? & --- & 0 & --- & --- & --- & --- & --- & sem pares suficientes \\
regiao\_confusa & ? & 0.396 & 71 & -0.075 & 0.0766 & --- & --- & --- & indicativo \\
regiao\_incerteza & ? & --- & 0 & --- & --- & --- & --- & --- & sem pares suficientes \\
sorteada & ? & --- & 0 & --- & --- & --- & --- & --- & sem pares suficientes \\
sorteado & ? & --- & 0 & --- & --- & --- & --- & --- & sem pares suficientes \\
uniforme & ? & --- & 0 & --- & --- & --- & --- & --- & sem pares suficientes \\
uniforme\_balanceado & ? & --- & 0 & --- & --- & --- & --- & --- & sem pares suficientes \\
\bottomrule
\end{tabular}
\end{adjustbox}

\end{table}

% Nota de comparacao_criterios. GERADA por flim_al/tabelas.py.
\footnotesize \texttt{random} é o piso --- sem critério nenhum. O pareamento casa split, decoder, orçamento e semente \textbf{dentro da mesma família de experimento}; a média mostrada vem dessas mesmas células, para que média e $\Delta$ falem do mesmo conjunto. \textbf{Toda técnica que rodou está aqui, inclusive as que perderam do sorteio} --- no experimento de encoder retreinado isso aconteceu com as quatro de imagem. Negrito só com p < 0,05. Para técnicas de REGIÃO, leia a coluna do controle de nível (\texttt{random\_region}): o $\Delta$ contra \texttt{random} soma seleção e geometria, e a geometria pesa ~2$\times$ a seleção neste projeto. \textbf{\texttt{flim\_3img} é o FLIM do artigo} --- as 3 imagens que os autores escolheram, sem seleção automática. A coluna \texttt{$\Delta$ vs FLIM} responde à pergunta central: aplicar AL melhora sobre isso? Ela só aparece em orçamento casado (K=3), porque o braço do artigo só existe com 3 imagens; comparar AL com 10 imagens contra FLIM com 3 mediria orçamento, não seleção. \textbf{Coluna toda vazia significa comparação IMPOSSÍVEL com os dados existentes, não ausente por descuido}: o braço do artigo usa os 31 markers REAIS do especialista, e os braços de AL sobre o pool completo usam sintéticos, porque marker real só existe para 31 imagens. \texttt{marker\_origem} entra no pareamento para que esse par não se forme --- ele mediria quem desenhou o traço e apresentaria como efeito da seleção. Para responder à pergunta, rode \texttt{scripts/varrer\_criterios.py --modo imagem}: marker real nos dois lados, com \texttt{oracle} (o passo 7 do Algoritmo 1) como referência.
\normalsize


Linhas com células vazias merecem explicação, porque a ausência ali não é descuido. O braço do trabalho de referência usa os trinta e um conjuntos de marcadores \emph{reais} dos especialistas, enquanto os braços de \ac{AL} sobre o conjunto completo usam marcadores sintéticos, pela razão simples de que marcador real só existe para essas trinta e uma imagens. A origem do marcador entra no pareamento justamente para impedir que esse par se forme: se formasse, a diferença mediria quem desenhou o traço e seria apresentada como efeito da seleção. Onde a coluna aparece vazia, a comparação é \textbf{impossível} com os dados existentes.

A Tabela~\ref{tab_custo} traz o custo computacional medido, que é parte do argumento a favor do \ac{FLIM} e precisa ser registrado ao lado da qualidade. Estimar os núcleos por agrupamento leva de 0,8 segundo com uma imagem a 48,1 segundos com oito. O tempo de avaliação cresce com o orçamento pela razão já discutida: o codificador cresce junto. Todas as linhas vêm de uma única campanha com uma semente, de modo que servem para ordem de grandeza e não para comparar critérios entre si.

\begin{table}[htb]
\centering
\caption{Qualidade por critério e orçamento nos sete decodificadores, com custo.}
\label{tab_custo}
% Fβ por critério e orçamento, nos sete decoders, com custo
% GERADO por flim_al/tabelas.py em 2026-10-08T13:52:40+00:00
% NAO EDITE A MAO: a proxima geracao sobrescreve, e o numero editado deixa de bater com evidencia/execucoes/.
% Proveniencia: k_por_modelo.proveniencia.json
\begin{adjustbox}{max width=\linewidth}
\begin{tabular}{llrrrrrrrrrrrr}
\toprule
critério & K & n & FLIM\_lm & FLIM\_pb & FLIM\_mb & FLIM\_at & FLIM\_ts & FLIM\_lt & FLIM\_ts* & melhor & $\Delta$ vs artigo & treino (s) & teste (s) \\
\midrule
flim\_paper & 1 & 1 & 0.443 & 0.530 & 0.341 & 0.420 & 0.329 & 0.294 & 0.362 & 0.530 (FLIM\_pb) & 0.000 & 0.8 & 85.4 \\
flim\_paper & 2 & 1 & 0.433 & 0.574 & 0.417 & 0.299 & 0.460 & 0.409 & 0.273 & 0.574 (FLIM\_pb) & 0.000 & 2.4 & 125.8 \\
flim\_paper & 3 & 1 & 0.133 & 0.623 & 0.213 & 0.125 & 0.240 & 0.242 & 0.243 & 0.623 (FLIM\_pb) & 0.000 & 4.3 & 160.6 \\
flim\_paper & 4 & 1 & 0.575 & 0.634 & 0.418 & 0.196 & 0.488 & 0.610 & 0.403 & 0.634 (FLIM\_pb) & 0.000 & 7.9 & 210.4 \\
flim\_paper & 5 & 1 & 0.628 & 0.659 & 0.415 & 0.364 & 0.470 & 0.535 & 0.378 & 0.659 (FLIM\_pb) & 0.000 & 11.9 & 211.6 \\
coreset & 1 & 1 & 0.661 & 0.641 & 0.341 & 0.176 & 0.323 & 0.277 & 0.334 & 0.661 (FLIM\_lm) & 0.005 & 0.9 & 97.7 \\
coreset & 2 & 1 & 0.657 & 0.491 & 0.217 & 0.209 & 0.167 & 0.136 & 0.293 & 0.657 (FLIM\_lm) & -0.099 & 3.2 & 143.0 \\
coreset & 3 & 1 & 0.323 & 0.594 & 0.199 & 0.167 & 0.104 & 0.179 & 0.240 & 0.594 (FLIM\_pb) & -0.001 & 4.3 & 182.1 \\
coreset & 4 & 1 & 0.604 & 0.561 & 0.329 & 0.249 & 0.308 & 0.317 & 0.326 & 0.604 (FLIM\_lm) & -0.090 & 9.3 & 246.2 \\
coreset & 5 & 1 & 0.460 & 0.683 & 0.462 & 0.372 & 0.498 & 0.539 & 0.385 & 0.683 (FLIM\_pb) & -0.007 & 14.0 & 257.2 \\
coreset & 8 & 1 & 0.682 & 0.617 & 0.379 & 0.288 & 0.441 & 0.487 & 0.300 & 0.682 (FLIM\_lm) & --- & 25.5 & 257.5 \\
entropy & 1 & 1 & 0.069 & 0.419 & 0.098 & 0.103 & 0.136 & 0.147 & 0.197 & 0.419 (FLIM\_pb) & -0.222 & 0.8 & 92.2 \\
entropy & 2 & 1 & 0.532 & 0.630 & 0.351 & 0.300 & 0.237 & 0.249 & 0.321 & 0.630 (FLIM\_pb) & -0.035 & 3.0 & 152.6 \\
entropy & 3 & 1 & 0.292 & 0.563 & 0.589 & 0.163 & 0.550 & 0.587 & 0.403 & 0.589 (FLIM\_mb) & 0.190 & 4.6 & 190.7 \\
entropy & 4 & 1 & 0.586 & 0.651 & 0.407 & 0.273 & 0.337 & 0.396 & 0.382 & 0.651 (FLIM\_pb) & -0.042 & 9.2 & 251.0 \\
entropy & 5 & 1 & 0.397 & 0.714 & 0.400 & 0.267 & 0.255 & 0.355 & 0.390 & 0.714 (FLIM\_pb) & -0.096 & 20.6 & 250.3 \\
entropy & 8 & 1 & 0.611 & 0.465 & 0.284 & 0.200 & 0.368 & 0.302 & 0.260 & 0.611 (FLIM\_lm) & --- & 25.0 & 252.5 \\
least\_confidence & 1 & 1 & 0.069 & 0.419 & 0.098 & 0.103 & 0.136 & 0.147 & 0.197 & 0.419 (FLIM\_pb) & -0.222 & 1.4 & 92.8 \\
least\_confidence & 2 & 1 & 0.532 & 0.630 & 0.351 & 0.300 & 0.237 & 0.249 & 0.321 & 0.630 (FLIM\_pb) & -0.035 & 2.6 & 146.8 \\
least\_confidence & 3 & 1 & 0.292 & 0.563 & 0.589 & 0.163 & 0.550 & 0.587 & 0.403 & 0.589 (FLIM\_mb) & 0.190 & 6.0 & 188.6 \\
least\_confidence & 4 & 1 & 0.648 & 0.639 & 0.237 & 0.287 & 0.194 & 0.232 & 0.315 & 0.648 (FLIM\_lm) & -0.110 & 9.1 & 250.3 \\
least\_confidence & 5 & 1 & 0.529 & 0.675 & 0.300 & 0.250 & 0.279 & 0.273 & 0.290 & 0.675 (FLIM\_pb) & -0.122 & 17.8 & 249.2 \\
least\_confidence & 8 & 1 & 0.626 & 0.495 & 0.277 & 0.201 & 0.337 & 0.300 & 0.251 & 0.626 (FLIM\_lm) & --- & 27.6 & 246.9 \\
medoide & 1 & 1 & 0.623 & 0.535 & 0.219 & 0.115 & 0.171 & 0.237 & 0.269 & 0.623 (FLIM\_lm) & -0.079 & 1.3 & 106.6 \\
medoide & 2 & 1 & 0.644 & 0.530 & 0.171 & 0.143 & 0.098 & 0.141 & 0.290 & 0.644 (FLIM\_lm) & -0.121 & 2.9 & 155.9 \\
medoide & 3 & 1 & 0.676 & 0.623 & 0.289 & 0.098 & 0.358 & 0.367 & 0.293 & 0.676 (FLIM\_lm) & 0.127 & 4.6 & 197.0 \\
medoide & 4 & 1 & 0.655 & 0.522 & 0.238 & 0.079 & 0.296 & 0.333 & 0.294 & 0.655 (FLIM\_lm) & -0.129 & 10.9 & 253.7 \\
medoide & 5 & 1 & 0.670 & 0.490 & 0.323 & 0.183 & 0.342 & 0.376 & 0.307 & 0.670 (FLIM\_lm) & -0.108 & 14.7 & 249.9 \\
medoide & 8 & 1 & 0.669 & 0.582 & 0.375 & 0.235 & 0.405 & 0.404 & 0.309 & 0.669 (FLIM\_lm) & --- & 27.4 & 247.9 \\
random & 1 & 1 & 0.701 & 0.695 & 0.452 & 0.142 & 0.289 & 0.298 & 0.444 & 0.701 (FLIM\_lm) & 0.043 & 0.9 & 83.6 \\
random & 2 & 1 & 0.611 & 0.690 & 0.305 & 0.346 & 0.321 & 0.297 & 0.320 & 0.690 (FLIM\_pb) & 0.004 & 2.5 & 130.3 \\
random & 3 & 1 & 0.560 & 0.581 & 0.421 & 0.231 & 0.498 & 0.491 & 0.362 & 0.581 (FLIM\_pb) & 0.190 & 6.5 & 173.4 \\
random & 4 & 1 & 0.397 & 0.561 & 0.397 & 0.260 & 0.407 & 0.471 & 0.362 & 0.561 (FLIM\_pb) & -0.067 & 7.5 & 210.7 \\
random & 5 & 1 & 0.632 & 0.575 & 0.346 & 0.315 & 0.369 & 0.365 & 0.364 & 0.632 (FLIM\_lm) & -0.069 & 12.1 & 215.0 \\
random & 8 & 1 & 0.681 & 0.607 & 0.489 & 0.508 & 0.566 & 0.551 & 0.442 & 0.681 (FLIM\_lm) & --- & 48.1 & 266.5 \\
\bottomrule
\end{tabular}
\end{adjustbox}

\end{table}

% Nota de k_por_modelo. GERADA por flim_al/tabelas.py.
\footnotesize Uma linha por (critério, K); cada coluna de decoder é o F$\beta$ naquele decoder. \texttt{flim\_paper} é o braço do artigo --- as imagens fixas que os autores escolheram, sem seleção automática. \textbf{Todos os braços usam o mesmo gerador de marker sintético}, inclusive o do artigo: marker real só existe para 31 imagens, e misturar as duas origens faria o $\Delta$ medir quem desenhou o traço em vez da seleção. A comparação é internamente válida e \textbf{não} se compara com números produzidos a partir de markers reais. \texttt{treino} é estimar os kernels por k-means --- o que o FLIM promete ser barato; \texttt{teste} é rodar os sete decoders no conjunto de avaliação. Ele cresce com K porque o encoder cresce: a arquitetura pede 200 kernels por camada, mas o k-means só produz tantos quantos os patches dos markers permitem --- medido, 54/51/48/48 com uma imagem contra 200/200/200/200 com oito. Com poucas anotações a rede não é só menos treinada, é menor. Tempo com \texttt{~} é estimado a partir da soma treino+avaliação, que era tudo que o registro guardava antes; sem \texttt{~} é medido. $\Delta$ é contra o artigo no MESMO K, e fica vazio em K acima de cinco porque o braço do artigo tem cinco imagens e não existe referência para comparar. Todas as linhas vêm de uma única campanha (\texttt{tabela\_k\_por\_modelo/semente0/pool40/teste60}): campanhas diferentes usam conjuntos de teste diferentes e não se misturam na mesma linha.
\normalsize


\section{Seleção de região: a margem existe}

A ausência de ganho na seleção de imagens admite duas leituras opostas. Pode ser que a escolha não importe, caso em que a técnica seria inaplicável por natureza. Ou pode ser que a escolha importe e os critérios não a capturem, caso em que o problema permanece aberto. Distinguir as duas exige medir o teto, e não apenas o desempenho.

O experimento de ganho marginal fixa a partição, as imagens de treino, os marcadores iniciais, o codificador de partida e o conjunto de validação, e varia uma única coisa: qual região recebe os aproximadamente trezentos pixels seguintes de anotação. Para cada região candidata o codificador é reestimado e o $F_\beta$ é medido novamente, de modo que a diferença entre candidatos isola o efeito da posição da anotação.

\begin{table}[htb]
\centering
\caption{Ganho marginal de uma anotação adicional no conjunto de parasitas.}
\label{tab_ganho_marginal}
% Ganho marginal de um marcador adicional — Parasitas, FLIM_lm, base funcional
% GERADO por flim_al/tabelas.py em 2026-10-06T16:49:57+00:00
% NAO EDITE A MAO: a proxima geracao sobrescreve, e o numero editado deixa de bater com evidencia/execucoes/.
% Proveniencia: ganho_marginal_schisto_FLIM_lm_funcional.proveniencia.json
\begin{tabular}{lrrrrrrr}
\hline
o que recebeu o marcador & n (sementes) & $\Delta$ F$\beta$ médio & IC95\% & pior & melhor & $\Delta$ vs sorteado & p \\
\hline
melhor região (teto amostrado) & 10 & 0.2191 & [0.0851, 0.3530] & 0.0300 & 0.5777 & 0.3089 & 0.0000 \\
argmax entropia (o AL) & 10 & 0.0843 & [-0.0577, 0.2264] & -0.1196 & 0.4676 & 0.1741 & 0.0157 \\
argmax least confidence & 10 & 0.1078 & [-0.0267, 0.2423] & -0.1168 & 0.4647 & 0.1976 & 0.0071 \\
candidato sorteado (média) & 10 & -0.0898 & [-0.1697, -0.0099] & -0.2672 & 0.1461 & --- & --- \\
pior região & 10 & -0.3490 & [-0.4764, -0.2216] & -0.6091 & -0.0433 & -0.2592 & 0.0000 \\
imagem nova inteira (orçamento MAIOR) & 10 & 0.0506 & [-0.0402, 0.1414] & -0.1427 & 0.3172 & 0.1404 & 0.0002 \\
sorteio uniforme (reponderado, sem GT) & 10 & -0.1050 & [-0.1789, -0.0311] & -0.3038 & 0.1034 & -0.0152 & 0.1930 \\
\hline
\end{tabular}

\end{table}

% Nota de ganho_marginal_schisto_FLIM_lm_funcional. GERADA por flim_al/tabelas.py.
\footnotesize $\Delta$ = F$\beta$ na \textbf{validação} com o marcador extra menos F$\beta$ sem ele. Dentro de cada semente, a partição, as imagens iniciais, os marcadores base, o encoder inicial e o conjunto de validação são idênticos --- a única variável é qual região recebeu os ~300 px adicionais. O treino é determinístico \textbf{dentro de um processo}, a partir dos mesmos arquivos de marcador (verificado: três repetições dão F$\beta$ idêntico até a décima casa). \textbf{Entre} execuções há ruído: a reexecução completa do schisto deu F$\beta$ idêntico em 538 de 540 registros, e nos 2 restantes o desvio máximo foi 0,0034 --- duas ordens de grandeza abaixo dos efeitos reportados. \textbf{n é o número de sementes}: os candidatos de uma mesma semente compartilham encoder e são correlacionados. \texttt{melhor região (teto)} escolhe pelo F$\beta$ da validação e \textbf{não é uma estratégia} --- mede o que existe para capturar, \textbf{entre os 22 candidatos examinados} de ~380 superpixels; não é teto absoluto. Linhas restritas a: \textbf{base funcional}. A separação por regime é post-hoc e forçada pela aritmética --- $\Delta$ a partir de F$\beta$=0 exato não pode ser negativo, então as duas populações não têm média comparável. As sementes compartilham o pool e o conjunto de validação: o IC95\% vale para \textbf{esta} validação sob sorteio das imagens de treino, e não generaliza para o dataset. O conjunto de teste não foi tocado por esta campanha.
\normalsize


O mesmo experimento nos outros dois conjuntos, nas Tabelas~\ref{tab_gm_brats} e \ref{tab_gm_conj}, reproduz o padrão com magnitudes diferentes, o que é o que permite afirmar que a margem não é particularidade de um domínio.

\begin{table}[htb]
\centering
\caption{Ganho marginal de uma anotação adicional no BraTS.}
\label{tab_gm_brats}
% Ganho marginal de um marcador adicional — BraTS, FLIM_lm, base funcional
% GERADO por flim_al/tabelas.py em 2026-10-07T13:57:17+00:00
% NAO EDITE A MAO: a proxima geracao sobrescreve, e o numero editado deixa de bater com evidencia/execucoes/.
% Proveniencia: ganho_marginal_brats_FLIM_lm_funcional.proveniencia.json
\begin{tabular}{lrrrrrrr}
\hline
o que recebeu o marcador & n (sementes) & $\Delta$ F$\beta$ médio & IC95\% & pior & melhor & $\Delta$ vs sorteado & p \\
\hline
melhor região (teto amostrado) & 7 & 0.0748 & [-0.0049, 0.1544] & 0.0080 & 0.2581 & 0.2766 & 0.0026 \\
argmax entropia (o AL) & 7 & -0.1134 & [-0.3468, 0.1199] & -0.6785 & 0.0532 & 0.0884 & 0.2847 \\
argmax least confidence & 7 & 0.0301 & [-0.0027, 0.0630] & -0.0344 & 0.0680 & 0.2320 & 0.0012 \\
candidato sorteado (média) & 7 & -0.2019 & [-0.2966, -0.1071] & -0.3571 & -0.0788 & --- & --- \\
pior região & 7 & -0.6972 & [-0.7849, -0.6095] & -0.7768 & -0.4949 & -0.4954 & 0.0001 \\
imagem nova inteira (orçamento MAIOR) & 7 & -0.2158 & [-0.4314, -0.0001] & -0.5292 & 0.0363 & -0.0139 & 0.8895 \\
sorteio uniforme (reponderado, sem GT) & 7 & -0.2224 & [-0.3593, -0.0856] & -0.4400 & -0.0381 & -0.0206 & 0.3773 \\
\hline
\end{tabular}

\end{table}

\begin{table}[htb]
\centering
\caption{Ganho marginal de uma anotação adicional na conjuntivite.}
\label{tab_gm_conj}
% Ganho marginal de um marcador adicional — Conjuntivite, FLIM_lm, base funcional
% GERADO por flim_al/tabelas.py em 2026-10-08T13:52:42+00:00
% NAO EDITE A MAO: a proxima geracao sobrescreve, e o numero editado deixa de bater com evidencia/execucoes/.
% Proveniencia: ganho_marginal_conjunctiva_FLIM_lm_funcional.proveniencia.json
\begin{adjustbox}{max width=\linewidth}
\begin{tabular}{lrrrrrrr}
\toprule
o que recebeu o marcador & n (sementes) & $\Delta$ F$\beta$ médio & IC95\% & pior & melhor & $\Delta$ vs sorteado & p \\
\midrule
melhor região (teto amostrado) & 10 & 0.1266 & [0.0303, 0.2229] & 0.0031 & 0.4514 & 0.1997 & 0.0001 \\
argmax entropia (o AL) & 10 & 0.0391 & [-0.0725, 0.1506] & -0.1786 & 0.4294 & 0.1121 & 0.0042 \\
argmax least confidence & 10 & 0.0268 & [-0.0771, 0.1307] & -0.2204 & 0.3591 & 0.0999 & 0.0070 \\
candidato sorteado (média) & 10 & -0.0731 & [-0.1790, 0.0329] & -0.3248 & 0.2256 & --- & --- \\
pior região & 10 & -0.5076 & [-0.6155, -0.3996] & -0.7023 & -0.2655 & -0.4345 & 0.0000 \\
imagem nova inteira (orçamento MAIOR) & 10 & -0.0290 & [-0.2288, 0.1708] & -0.6899 & 0.3608 & 0.0441 & 0.5379 \\
sorteio uniforme (reponderado, sem GT) & 10 & -0.0758 & [-0.1872, 0.0355] & -0.3512 & 0.2254 & -0.0028 & 0.6053 \\
\bottomrule
\end{tabular}
\end{adjustbox}

\end{table}

% Nota de ganho_marginal_conjunctiva_FLIM_lm_funcional. GERADA por flim_al/tabelas.py.
\footnotesize $\Delta$ = F$\beta$ na \textbf{validação} com o marcador extra menos F$\beta$ sem ele. Dentro de cada semente, a partição, as imagens iniciais, os marcadores base, o encoder inicial e o conjunto de validação são idênticos --- a única variável é qual região recebeu os ~300 px adicionais. O treino é determinístico \textbf{dentro de um processo}, a partir dos mesmos arquivos de marcador (verificado: três repetições dão F$\beta$ idêntico até a décima casa). \textbf{Entre} execuções há ruído: a reexecução completa do schisto deu F$\beta$ idêntico em 538 de 540 registros, e nos 2 restantes o desvio máximo foi 0,0034 --- duas ordens de grandeza abaixo dos efeitos reportados. \textbf{n é o número de sementes}: os candidatos de uma mesma semente compartilham encoder e são correlacionados. \texttt{melhor região (teto)} escolhe pelo F$\beta$ da validação e \textbf{não é uma estratégia} --- mede o que existe para capturar, \textbf{entre os 22 candidatos examinados} de ~380 superpixels; não é teto absoluto. Linhas restritas a: \textbf{base funcional}. A separação por regime é post-hoc e forçada pela aritmética --- $\Delta$ a partir de F$\beta$=0 exato não pode ser negativo, então as duas populações não têm média comparável. As sementes compartilham o pool e o conjunto de validação: o IC95\% vale para \textbf{esta} validação sob sorteio das imagens de treino, e não generaliza para o dataset. O conjunto de teste não foi tocado por esta campanha.
\normalsize


A linha de imagem nova inteira merece comentário, porque ela responde a uma objeção natural. Poderia alguém argumentar que, em vez de escolher melhor onde anotar numa imagem já escolhida, bastaria anotar uma imagem a mais. A linha mede exatamente isso, e com orçamento \emph{maior} que o das demais: ela rende menos que a melhor região. O ganho não está em anotar mais, está em anotar no lugar certo, e é esse contraste que justifica mudar a granularidade da pergunta de imagem para região.

A dispersão entre candidatos é grande. A melhor região examinada supera a região sorteada por margens entre 0,1351 e 0,3089 de $F_\beta$, com significância nas seis células formadas por conjunto de dados e decodificador, todas com valor p igual ou inferior a 0,0031. A pior região, por sua vez, fica entre 0,2075 e 0,4954 abaixo do sorteio, também nas seis células.

O segundo número merece tanta atenção quanto o primeiro. O risco é da mesma ordem de grandeza do prêmio, o que significa que, no \ac{FLIM}, \textbf{anotar mais não é monotonicamente bom}. Em uma rede treinada por retropropagação, acrescentar um exemplo rotulado correto raramente piora o modelo; aqui, uma única anotação adicional mal posicionada pode custar meio ponto de $F_\beta$. É essa assimetria que dá sentido à pergunta sobre onde anotar: se toda anotação ajudasse, bastaria anotar mais.

Dois esclarecimentos são necessários. A linha identificada como melhor região escolhe o candidato pelo $F_\beta$ de validação e portanto não constitui uma estratégia aplicável: ela mede o que existe para ser capturado, e não um método. E o valor medido é um máximo sobre os candidatos examinados, não um teto absoluto, já que apenas parte dos superpixels de cada imagem foi avaliada.

Ainda assim, o resultado responde à ambiguidade. Existe margem considerável na escolha da região, e ela é da ordem de grandeza que separa um modelo mediano de um modelo bom. A pergunta passa a ser quanto dessa margem os critérios alcançam.

\section{O custo do Active Learning}
\label{sec_custo_al}

Esta seção reúne as duas metades do resultado central. De um lado, quanto vale o clique no lugar certo, medido pelo teto amostrado. De outro, quanto o critério de \ac{AL} efetivamente entrega quando escolhe sozinho, sem olhar resultado nenhum. A diferença entre as duas é o que se deixa na mesa, e é ela que esta dissertação chama de \emph{custo do Active Learning}.

As duas colunas vêm da mesma campanha, da mesma semente, com a mesma base e a mesma validação, de modo que a subtração é pareada e não mistura condições.

\begin{table}[htb]
\centering
\caption{A margem disponível, o que o critério captura, e a lacuna entre as duas.}
\label{tab_custo_al}
% O custo do Active Learning: a margem que existe e a que se captura
% GERADO por flim_al/tabelas.py em 2026-10-08T14:01:07+00:00
% NAO EDITE A MAO: a proxima geracao sobrescreve, e o numero editado deixa de bater com evidencia/execucoes/.
% Proveniencia: custo_do_al.proveniencia.json
\begin{adjustbox}{max width=\linewidth}
\begin{tabular}{llrrrrrrrr}
\toprule
conjunto & decodificador & n & máximo amostrado & o critério & lacuna & p & q (BH) & sinais & captura \\
\midrule
Parasitas & FLIM\_lm & 10 & 0.3089 & 0.1741 & \textbf{0.1347} & 0.0230 & 0.0345 & 0.3438 & 56\% \\
Parasitas & FLIM\_pb & 10 & 0.1351 & 0.0096 & \textbf{0.1255} & 0.0006 & 0.0039 & 0.0020 & 7\% \\
BraTS & FLIM\_lm & 7 & 0.2766 & 0.0884 & 0.1882 & 0.1233 & 0.1233 & 0.1250 & 32\% \\
BraTS & FLIM\_pb & 8 & 0.2384 & -0.0283 & 0.2667 & 0.0495 & 0.0594 & 0.2891 & -12\% \\
Conjuntivite & FLIM\_lm & 10 & 0.1997 & 0.1121 & \textbf{0.0876} & 0.0038 & 0.0115 & 0.0020 & 56\% \\
Conjuntivite & FLIM\_pb & 10 & 0.1436 & 0.0598 & \textbf{0.0839} & 0.0117 & 0.0234 & 0.0215 & 42\% \\
\bottomrule
\end{tabular}
\end{adjustbox}

\end{table}

% Nota de custo_do_al. GERADA por flim_al/tabelas.py.
\footnotesize Todas as colunas são diferenças de F$\beta$ contra o \textbf{candidato sorteado}, medidas na validação, com a partição, as imagens, os marcadores base e o codificador inicial idênticos dentro de cada semente. \texttt{teto amostrado} é o melhor candidato entre os examinados, escolhido \textbf{pelo F$\beta$ da validação}: não é estratégia, e mede o que existe para capturar, não um teto absoluto. \texttt{o critério} é o que o Active Learning de fato escolhe, sem olhar resultado. A \textbf{lacuna} entre os dois é o custo de o critério errar o lugar do clique, e é a coluna com teste: t pareado por semente, com n = sementes, porque os candidatos de uma mesma semente compartilham codificador. \texttt{captura} é razão entre duas médias, fica instável com denominador pequeno e serve de ilustração, não de estatística. Linhas restritas à \textbf{base funcional}; a estratificação por regime é post-hoc e forçada pela aritmética, porque $\Delta$ a partir de F$\beta$=0 exato não pode ser negativo. As sementes compartilham pool e validação, então o IC vale para esta validação sob sorteio das imagens de treino e não generaliza para o dataset. O conjunto de teste não foi tocado por esta campanha.
\normalsize


Três leituras, em ordem de importância.

\textbf{A margem é grande e consistente.} De 0,1351 a 0,3089 de $F_\beta$, positiva nas seis células, com significância em todas. Esse é o valor do clique no lugar certo, e é a razão de a pergunta sobre posição de anotação valer a pena. Para dar escala: no mesmo conjunto de dados, a diferença entre o melhor e o pior dos sete decodificadores avaliados é da ordem de 0,2, de modo que escolher onde anotar pesa tanto quanto escolher a arquitetura do decodificador.

\textbf{O critério captura parte dela, e a parte não capturada é significativa.} A lacuna vai de 0,0839 a 0,2667 e atinge significância em quatro das seis células. Em termos de fração capturada, o critério recupera cerca de metade da margem nos casos favoráveis e praticamente nada nos desfavoráveis. A coluna de captura é apresentada como ilustração e não como estatística, porque é razão entre duas médias e fica instável quando o denominador é pequeno; quem tem teste é a lacuna.

\textbf{Há um caso em que o critério vai bem, e ele precisa ser dito.} No BraTS com o decodificador \texttt{FLIM\_lm}, o \textit{least confidence} entrega 0,2320 contra um teto de 0,2766, isto é, captura 84\% da margem, com valor p de 0,0012. Omitir esse caso tornaria a conclusão mais limpa e menos verdadeira. A leitura correta não é que os critérios de incerteza sejam inúteis, e sim que o seu desempenho \textbf{depende do decodificador e do domínio}. O mesmo escore que ordena bem as regiões no \texttt{FLIM\_lm} do conjunto de parasitas, com correlação de 0,427 e valor p de 0,0001, não as ordena no \texttt{FLIM\_pb}, onde a correlação é de $-0{,}045$ com valor p de 0,6400.

A formulação que os números autorizam é, portanto:

\begin{quote}
Sob orçamento de anotação fixo, existe uma margem de 0,1351 a 0,3089 de $F_\beta$ entre a melhor região examinada e uma região sorteada. Os critérios de \ac{AL} capturam parte dela, e a parte não capturada, de 0,0839 a 0,2667, é o custo de o critério escolher o lugar errado.
\end{quote}

Note o que a formulação \emph{não} diz. Ela não diz que o \ac{AL} melhora o $F_\beta$ em 0,1 a 0,3. Os 0,1 a 0,3 são o teto, medido com acesso ao resultado da validação, e nenhum método os entrega. Confundir as duas coisas transformaria o resultado mais sólido do trabalho em uma afirmação falsa.

\section{Para onde os critérios apontam}

A resposta está em medir o comportamento dos critérios, e não apenas o seu desempenho. Como escolher uma região exige apenas pontuar e ordenar, essa medição dispensa reestimar o codificador para cada candidato e pode ser feita sobre dez repetições nos três conjuntos de dados a custo baixo.

A fração de objeto de uma região descreve onde ela está: valor próximo de zero indica fundo puro, valor próximo de um indica interior do objeto, e valores intermediários indicam que a região atravessa a fronteira. Essa distinção importa porque a comparação com os marcadores reais dos especialistas, apresentada na Tabela \ref{tab_artigo_regiao}, mede que a anotação de borda vale entre 0,126 e 0,166 de $F_\beta$ sobre a anotação em qualquer outro ponto interno ao objeto.

\begin{table}[htb]
\centering
\caption{Anotação dos especialistas comparada a cliques realocados.}
\label{tab_artigo_regiao}
% Tabela do artigo × cliques realocados por AL × realocados ao acaso
% GERADO por flim_al/tabelas.py em 2026-10-08T14:01:04+00:00
% NAO EDITE A MAO: a proxima geracao sobrescreve, e o numero editado deixa de bater com evidencia/execucoes/.
% Proveniencia: artigo_vs_regiao.proveniencia.json
\begin{adjustbox}{max width=\linewidth}
\begin{tabular}{lrrrrrrrrrrr}
\toprule
decoder & n & F$\beta$ artigo & F$\beta$ AL região & F$\beta$ aleatório & $\Delta$ AL-artigo & p & $\Delta$ aleat-artigo & p  & $\Delta$ AL-aleat & p   & F$\beta$ publicado (A) \\
\midrule
FLIM\_lm & 6 & 0.744 & 0.622 & 0.669 & -0.122 & 0.0370 & -0.075 & 0.0480 & -0.047 & 0.3088 & 0.860 \\
FLIM\_pb & 6 & 0.720 & 0.494 & 0.537 & -0.225 & 0.0140 & -0.182 & 0.0240 & -0.043 & 0.4223 & 0.857 \\
FLIM\_mb & 6 & 0.718 & 0.502 & 0.544 & -0.217 & 0.0106 & -0.174 & 0.0121 & -0.042 & 0.3997 & 0.843 \\
FLIM\_at & 6 & 0.518 & 0.303 & 0.354 & -0.214 & 0.0051 & -0.164 & 0.0015 & -0.050 & 0.1178 & 0.740 \\
FLIM\_ts & 6 & 0.591 & 0.466 & 0.491 & -0.124 & 0.0031 & -0.099 & 0.0000 & -0.025 & 0.2922 & 0.747 \\
FLIM\_lt & 6 & 0.649 & 0.549 & 0.575 & -0.100 & 0.0043 & -0.075 & 0.0238 & -0.026 & 0.2467 & 0.810 \\
FLIM\_ts* & 6 & 0.506 & 0.349 & 0.391 & -0.158 & 0.0029 & -0.115 & 0.0005 & -0.042 & 0.1427 & --- \\
\bottomrule
\end{tabular}
\end{adjustbox}

\end{table}

% Nota de artigo_vs_regiao. GERADA por flim_al/tabelas.py.
\footnotesize Mesmas imagens --- as que os especialistas A e B anotaram ---, mesmo número de cliques e mesmo traço de fundo (copiado, não regerado) nos três braços. Só muda para onde vão os cliques de objeto. \textbf{Não há seleção de imagem: é seleção de região.} As três diferenças decompõem o efeito --- \texttt{aleat-artigo} isola o ato de tirar o clique da borda, e \texttt{AL-aleat} isola a escolha do Active Learning. Sem a terceira, um $\Delta$ negativo na primeira seria creditado ao AL quando pode ser todo da segunda. Cada célula é um par (usuário, split), n = 6, pareado. \texttt{F$\beta$ publicado (A)} vem da Tabela III do artigo e \textbf{não} é comparável com a coluna reproduzida: o artigo aplica Dynamic Trees, cujo binário não executa nesta máquina. Avaliação em 250 imagens de Z$_1$\textbackslash{}T, as mesmas em todos os braços.
\normalsize


A medição do comportamento dos critérios produz três constatações.

A primeira é que os critérios de incerteza encontram o objeto. Regiões de interior puro representam de 2\% a 8\% do conjunto de candidatos, e entropia e \textit{least confidence} selecionam interior em 40\% a 100\% das escolhas, o que é um enriquecimento expressivo sobre a taxa de base.

A segunda é que eles erram a geometria. As escolhas atravessam a fronteira em 0\% a 20\% dos casos, embora seja justamente a travessia que a comparação com os especialistas mostra ter valor.

A terceira é a mais consequente: regiões que atravessam a fronteira representam apenas 2\% a 12\% do conjunto de candidatos. A anotação valiosa quase não está disponível para ser escolhida. Isso não é uma falha dos critérios, e sim do procedimento que gera a lista. O SLIC é construído para que as fronteiras dos superpixels coincidam com as bordas presentes na imagem, de modo que seus segmentos tendem a ficar inteiramente dentro ou inteiramente fora do objeto. A propriedade que faz dele um bom algoritmo de superpixels é a mesma que o torna inadequado como gerador de candidatos para anotação de borda.

A Figura~\ref{fig_gargalo} torna visível a terceira constatação, que é a mais consequente das três. Os superpixels do \ac{SLIC} tendem a ficar inteiramente de um lado da fronteira; os discos centrados no contorno a atravessam por construção.

\begin{figure}[htb]
\centering
\includegraphics[width=\linewidth]{gerador_candidatos.png}
\caption[O gargalo na geração de candidatos]{Os dois geradores de candidatos sobre a mesma imagem. Em vermelho, os candidatos que atravessam a fronteira verdadeira do objeto. O painel da direita usa predição ideal e ilustra a geometria do gerador proposto; com a predição real do codificador os discos caem no interior, e o ganho medido ocorre apesar disso.}
\label{fig_gargalo}
\end{figure}

A Tabela~\ref{tab_mecanismo} completa o quadro medindo o que \emph{prevê} o ganho de anotar uma região. A entropia da região correlaciona com o ganho no decodificador \texttt{FLIM\_lm} do conjunto de parasitas, com coeficiente de 0,427 e valor p de 0,0001, mas não no \texttt{FLIM\_pb}, onde o coeficiente é de $-0{,}045$ com valor p de 0,6400. O escore não é inútil; ele é dependente do decodificador.

\begin{table}[htb]
\centering
\caption{O que prevê o ganho de anotar uma região.}
\label{tab_mecanismo}
% O que prevê o ganho de anotar uma região — Parasitas, FLIM_lm
% GERADO por flim_al/tabelas.py em 2026-10-06T16:49:57+00:00
% NAO EDITE A MAO: a proxima geracao sobrescreve, e o numero editado deixa de bater com evidencia/execucoes/.
% Proveniencia: ganho_mecanismo_schisto_FLIM_lm_funcional.proveniencia.json
\begin{tabular}{lrrrr}
\hline
variável & $\rho$ de Spearman com $\Delta$ F$\beta$ & IC95\% & n (sementes) & p ($\rho$ $\neq$ 0) \\
\hline
entropia da região & 0.427 & [0.292, 0.562] & 10 & 0.0001 \\
least confidence & 0.427 & [0.292, 0.562] & 10 & 0.0001 \\
fração de foreground & 0.308 & [0.221, 0.395] & 10 & 0.0000 \\
$\Delta$ kernels do encoder & -0.299 & [-0.429, -0.170] & 10 & 0.0005 \\
($\Delta$ F$\beta$ médio ao anotar região de objeto) & -0.0740 & [-0.1661, 0.0181] & 10 & --- \\
($\Delta$ F$\beta$ médio ao anotar região de fundo) & -0.1077 & [-0.1816, -0.0339] & 10 & --- \\
\hline
\end{tabular}

\end{table}

% Nota de ganho_mecanismo_schisto_FLIM_lm_funcional. GERADA por flim_al/tabelas.py.
\footnotesize $\rho$ calculado \textbf{dentro} de cada semente, sobre os candidatos sorteados daquela semente; são os $\rho$ por semente que entram no teste, com n = sementes. Agregar candidatos de sementes diferentes num $\rho$ único trataria observações correlacionadas como independentes. \texttt{entropia da região} é o escore que o Active Learning usa para decidir --- um $\rho$ indistinguível de zero significa que o escore não ordena as regiões por utilidade, e é isso que explicaria todos os resultados negativos das campanhas anteriores. \texttt{$\Delta$ kernels} mede capacidade, não posição: é quantos filtros a mais o encoder extraiu com o marcador extra.
\normalsize


A linha de variação do número de núcleos merece atenção separada. Ela tem correlação \emph{negativa} com o ganho, de $-0{,}299$ com valor p de 0,0005, isto é, acrescentar mais núcleos ao codificador tende a piorar o resultado. Essa linha mede capacidade, não posição, e foi ela que serviu de base ao mecanismo declarado antes da execução da intervenção descrita adiante. Correlação de núcleos com ganho não é, contudo, demonstração de causalidade, e o texto não a trata como tal.

Os critérios de diversidade, avaliados no nível de região pela primeira vez neste trabalho, comportam-se de forma ainda menos favorável. CoreSet localiza o objeto apenas no conjunto de conjuntivite, onde os objetos são grandes, e nos demais seleciona fundo em todas as escolhas, sem qualquer enriquecimento sobre a taxa de base. O medoide nunca localiza o objeto. A explicação é o desequilíbrio do conjunto de candidatos: a diversidade é calculada sobre a distribuição de aparência das regiões, e quando o objeto ocupa cerca de 3\% da imagem ele não participa dessa distribuição de forma relevante. Nos dois conjuntos de objetos pequenos, CoreSet e medoide, que são critérios opostos, produzem a mesma resposta.

\section{Posição ou composição? O controle que separa as duas}

Uma objeção precisa ser respondida antes de atribuir o efeito à posição da anotação. Escolher uma região incerta não escolhe apenas um \emph{lugar}: escolhe junto uma \emph{composição} de classes. Regiões incertas concentram muito mais objeto do que a média, cerca de 73\% contra 16\%, proporção medida antes da execução. Se o ganho viesse de anotar mais objeto, e não de anotar na posição certa, a conclusão do capítulo mudaria.

O experimento que separa as duas coisas fixa o orçamento em pixels anotados por imagem, usa as mesmas imagens em todos os braços, e acrescenta um braço de controle: anotação sorteada, mas com a \emph{mesma proporção} entre objeto e fundo que a seleção por incerteza produz. Se o ganho fosse de composição, esse braço o reproduziria.

\begin{table}[htb]
\centering
\caption{Onde anotar, com orçamento fixo em pixels por imagem.}
\label{tab_onde_marcar}
% Onde anotar, com orçamento fixo em pixels: Fβ e as duas diferenças
% GERADO por flim_al/tabelas.py em 2026-10-08T13:52:40+00:00
% NAO EDITE A MAO: a proxima geracao sobrescreve, e o numero editado deixa de bater com evidencia/execucoes/.
% Proveniencia: onde_marcar.proveniencia.json
\begin{adjustbox}{max width=\linewidth}
\begin{tabular}{lllrrrrrr}
\toprule
braço & px/imagem & decoder & n & F$\beta$ médio & $\Delta$ vs uniforme & p & $\Delta$ vs balanceado & p  \\
\midrule
regiao\_incerteza & 2000 & FLIM\_lm & 5 & 0.247 & -0.215 & 0.022 & -0.256 & 0.013 \\
regiao\_incerteza & 2000 & FLIM\_pb & 5 & 0.476 & -0.167 & 0.019 & 0.209 & 0.051 \\
regiao\_incerteza & 4000 & FLIM\_lm & 5 & 0.055 & -0.283 & 0.022 & -0.054 & 0.263 \\
regiao\_incerteza & 4000 & FLIM\_pb & 5 & 0.488 & -0.142 & 0.023 & 0.305 & 0.037 \\
regiao\_incerteza & 8000 & FLIM\_lm & 5 & 0.000 & -0.011 & 0.264 & 0.000 & 1.000 \\
regiao\_incerteza & 8000 & FLIM\_pb & 5 & 0.410 & -0.189 & 0.129 & 0.239 & 0.121 \\
regiao\_aleatoria & 2000 & FLIM\_lm & 5 & 0.437 & -0.024 & 0.566 & -0.066 & 0.489 \\
regiao\_aleatoria & 2000 & FLIM\_pb & 5 & 0.599 & -0.044 & 0.048 & 0.331 & 0.002 \\
regiao\_aleatoria & 4000 & FLIM\_lm & 5 & 0.387 & 0.049 & 0.412 & 0.278 & 0.037 \\
regiao\_aleatoria & 4000 & FLIM\_pb & 5 & 0.622 & -0.008 & 0.688 & 0.439 & 0.006 \\
regiao\_aleatoria & 8000 & FLIM\_lm & 5 & 0.070 & 0.059 & 0.303 & 0.070 & 0.229 \\
regiao\_aleatoria & 8000 & FLIM\_pb & 5 & 0.602 & 0.003 & 0.974 & 0.431 & 0.003 \\
regiao\_borda & 2000 & FLIM\_lm & 5 & 0.501 & 0.039 & 0.389 & -0.002 & 0.968 \\
regiao\_borda & 2000 & FLIM\_pb & 5 & 0.620 & -0.024 & 0.247 & 0.352 & 0.002 \\
regiao\_borda & 4000 & FLIM\_lm & 5 & 0.391 & 0.052 & 0.490 & 0.281 & 0.003 \\
regiao\_borda & 4000 & FLIM\_pb & 5 & 0.631 & 0.001 & 0.960 & 0.448 & 0.008 \\
regiao\_borda & 8000 & FLIM\_lm & 5 & 0.056 & 0.045 & 0.375 & 0.056 & 0.352 \\
regiao\_borda & 8000 & FLIM\_pb & 5 & 0.477 & -0.122 & 0.402 & 0.306 & 0.122 \\
uniforme\_balanceado & 2000 & FLIM\_lm & 5 & 0.503 & 0.042 & 0.606 & --- & --- \\
uniforme\_balanceado & 2000 & FLIM\_pb & 5 & 0.268 & -0.376 & 0.001 & --- & --- \\
uniforme\_balanceado & 4000 & FLIM\_lm & 5 & 0.109 & -0.229 & 0.111 & --- & --- \\
uniforme\_balanceado & 4000 & FLIM\_pb & 5 & 0.183 & -0.447 & 0.009 & --- & --- \\
uniforme\_balanceado & 8000 & FLIM\_lm & 5 & 0.000 & -0.011 & 0.264 & --- & --- \\
uniforme\_balanceado & 8000 & FLIM\_pb & 5 & 0.171 & -0.428 & 0.008 & --- & --- \\
\bottomrule
\end{tabular}
\end{adjustbox}

\end{table}

% Nota de onde_marcar. GERADA por flim_al/tabelas.py.
\footnotesize Orçamento em \textbf{pixels anotados por imagem}, não em imagens: o que limita o FLIM é a quantidade de pixels marcados --- a mesma imagem rende 54 kernels por camada com traço esparso e 200 com traço denso. Todos os braços gastam exatamente o mesmo número de pixels, nas MESMAS imagens (as do artigo); o que muda é onde eles caem. \texttt{uniforme} é o FLIM de hoje. \texttt{uniforme\_balanceado} tem a mesma proporção objeto/fundo que a incerteza produz, com os pixels sorteados --- é ele que separa *onde se anotou* de *quanto de cada classe se anotou*, porque escolher região incerta escolhe junto muito mais foreground (~73\% contra ~16\%, medido antes de rodar). \texttt{regiao\_borda} usa o gabarito e é teto, não método. Diferenças pareadas por semente; p de t pareado bicaudal. Campanha \texttt{onde\_marcar/k3/teste60/seg300}.
\normalsize


O resultado é desfavorável à seleção por incerteza em todas as seis células, com diferenças contra o braço uniforme entre $-0{,}011$ e $-0{,}283$. E o braço balanceado, que isola a composição, não reproduz o comportamento do braço de incerteza: o que separa os braços é o balanceamento de classe, e não a posição escolhida pelo escore. O braço de borda, que consulta o gabarito e portanto é teto e não método, fica dentro do ruído do uniforme.

A leitura conjunta com a Seção~\ref{sec_custo_al} é a seguinte. Existe margem na posição, e ela é grande. Mas o escore de incerteza, aplicado sobre superpixels, não a captura e ainda introduz um viés de composição que atrapalha. É o terceiro elo do argumento que leva ao gerador.

A Tabela~\ref{tab_regiao_vs_imagem} fecha a comparação por decodificador, e serve de lembrete de que os sete decodificadores de uma mesma célula compartilham o codificador: eles não são observações independentes entre si, e os valores p das suas linhas não se somam.

\begin{table}[htb]
\centering
\caption{Seleção por região contra o controle, por decodificador.}
\label{tab_regiao_vs_imagem}
% Active Learning por região: braços e controle, por decoder
% GERADO por flim_al/tabelas.py em 2026-10-08T14:01:02+00:00
% NAO EDITE A MAO: a proxima geracao sobrescreve, e o numero editado deixa de bater com evidencia/execucoes/.
% Proveniencia: regiao_vs_imagem.proveniencia.json
\begin{tabular}{lrrrrrr}
\toprule
decoder & al (média) & random (média) & n células & $\Delta$ & p & veredito \\
\midrule
decoder\_2 & 0.532 & 0.461 & 9 & 0.036 & 0.5113 & sem suporte \\
decoder\_3 & 0.395 & 0.316 & 9 & 0.033 & 0.6159 & sem suporte \\
decoder\_attention & 0.217 & 0.176 & 9 & 0.091 & 0.1743 & sem suporte \\
hybrid\_decoder & 0.394 & 0.384 & 9 & 0.079 & 0.1107 & sem suporte \\
labeled\_marker & 0.512 & 0.524 & 9 & -0.030 & 0.2452 & sem suporte \\
vanilla\_adaptive\_decoder & 0.429 & 0.413 & 9 & 0.030 & 0.4782 & sem suporte \\
vanilla\_adaptive\_decoder\_wt & 0.195 & 0.170 & 9 & \textbf{0.112} & 0.0471 & suportado \\
\bottomrule
\end{tabular}

\end{table}

% Nota de regiao_vs_imagem. GERADA por flim_al/tabelas.py.
\footnotesize $\Delta$ = al - random, teste \textbf{pareado}, pareado por semente. Só entram células em que os dois braços existem no \textbf{mesmo orçamento} --- o braço do artigo para em 3-4 imagens e o de AL vai a 9, e comparar orçamentos diferentes mediria duas coisas ao mesmo tempo. Negrito só com p < 0,05. Cada linha é um decoder; eles compartilham o encoder dentro de cada célula, então \textbf{não são observações independentes entre si} --- não some os p.
\normalsize


\section{Aprendizado contínuo sem retropropagação}

O esquecimento catastrófico, tal como descrito na Seção \ref{sec_cl}, é atribuído ao deslocamento dos pesos durante o ajuste por gradiente, e os métodos que o combatem atuam sobre esse deslocamento. Como o \ac{FLIM} não possui otimizador, nenhum desses métodos se aplica diretamente. O fenômeno, contudo, existe, e os resultados da seção anterior o quantificam: uma anotação adicional pode reduzir o $F_\beta$ em até meio ponto.

A inspeção do procedimento de estimação dos filtros localiza a origem. Os candidatos a filtro extraídos de todas as imagens e de todos os marcadores são reunidos em um único conjunto, que então é reduzido ao número de canais da camada por um agrupamento. Acrescentar marcadores acrescenta candidatos, o agrupamento é refeito sobre o conjunto ampliado, e os centróides que estavam em uso se deslocam. O banco de filtros é compartilhado por todas as imagens, inclusive as não anotadas, de modo que a perturbação é global e não local.

O mecanismo de retenção proposto no Capítulo \ref{cap_3} interpola o banco: cada centróide novo é ancorado no centróide anterior mais próximo, com um parâmetro de plasticidade que vai de retenção total a reestimação completa. A verificação de corretude confirma as duas extremidades, reproduzindo o comportamento original em um extremo e impedindo a entrada de qualquer direção nova no outro.

A avaliação, porém, não encontrou valor de plasticidade que convertesse o mecanismo em ganho mensurável. O teste primário, declarado antes da execução, comparou plasticidade intermediária contra reestimação completa e não atingiu significância, com sinais inconsistentes entre os tipos de região examinados. A retenção total chegou a parecer vantajosa em uma análise exploratória, mas o resultado não sobreviveu ao aumento do número de repetições nem ao teste em imagens inéditas, e o critério de falseamento declarado no pré-registro foi atingido: a qualidade final sob retenção total é menor, o que contraria a explicação proposta de que congelar o banco evitaria dano sem custo.

\section{Esforço de anotação}

O eixo da segmentação interativa é diferente do $F_\beta$ a orçamento fixo. A métrica usual é o número de cliques necessários para que a qualidade ultrapasse um alvo, com um teto de interações. Um método pode não alterar o $F_\beta$ a orçamento fixo e ainda assim reduzir o esforço de interação.

\begin{table}[htb]
\centering
\caption{Cliques até a qualidade alvo, por plasticidade do banco de filtros.}
\label{tab_noc}
% Cliques até a IoU alvo, por plasticidade do banco de filtros
% GERADO por flim_al/tabelas.py em 2026-10-08T14:01:07+00:00
% NAO EDITE A MAO: a proxima geracao sobrescreve, e o numero editado deixa de bater com evidencia/execucoes/.
% Proveniencia: noc_noc.proveniencia.json
\begin{adjustbox}{max width=\linewidth}
\begin{tabular}{lrrrrrrrrr}
\toprule
dataset & $\alpha$ & n & NoC médio & NoC mediano & taxa de sucesso & IoU final & $\Delta$ NoC vs $\alpha$=1 & p & McNemar \\
\midrule
Parasitas & 0.00 & 38 & 11.03 & 12.0 & 13\% & 0.5267 & -0.53 & 0.1031 & 3/0 (p=0.250) \\
Parasitas & 0.50 & 22 & 11.09 & 12.0 & 14\% & 0.5490 & -0.55 & 0.2876 & 2/0 (p=0.500) \\
Parasitas & 1.00 & 38 & 11.55 & 12.0 & 5\% & 0.5501 & --- & --- & --- \\
BraTS & 0.00 & 40 & 7.42 & 12.0 & 48\% & 0.4096 & -0.35 & 0.5829 & 4/3 (p=1.000) \\
BraTS & 0.50 & 20 & 6.00 & 3.0 & 60\% & 0.5273 & 0.05 & 0.7715 & 1/1 (p=1.000) \\
BraTS & 1.00 & 40 & 7.78 & 12.0 & 45\% & 0.4064 & --- & --- & --- \\
Conjuntivite & 0.00 & 1 & 12.00 & 12.0 & 0\% & 0.0000 & --- & --- & 0/1 (p=1.000) \\
Conjuntivite & 0.50 & 1 & 6.00 & 6.0 & 100\% & 0.7617 & --- & --- & 0/0 (p=1.000) \\
Conjuntivite & 1.00 & 1 & 6.00 & 6.0 & 100\% & 0.7764 & --- & --- & --- \\
\bottomrule
\end{tabular}
\end{adjustbox}

\end{table}

% Nota de noc_noc. GERADA por flim_al/tabelas.py.
\footnotesize NoC@X é o número de cliques até a IoU passar do alvo, com teto de 12; é o eixo de RITM (arXiv:2102.06583) e SimpleClick (arXiv:2210.11006), e \textbf{não} se compara com o NoC publicado deles, que usa outros conjuntos e outra definição de alvo. Imagem que não atinge o alvo entra com o teto, e por isso a \textbf{taxa de sucesso} aparece em toda linha: um braço pode mostrar NoC baixo por desistir mais cedo. \texttt{$\alpha$=1} é o FLIM puro, que refaz o banco de filtros a cada clique; \texttt{$\alpha$=0} congela o banco, e o clique chega ao modelo só pelos rótulos que o decoder \texttt{labeled\_marker} lê. Os testes são \textbf{separados por dataset}: agregar Schisto e BraTS, com 7\% e 30\% de sucesso, inflou um achado exploratório a p=0,0265 que separado não passa de p=0,10. O clique é simulado do ground truth pelo protocolo de Xu et al. (2016), com usuário que acerta sempre, então a tabela mede \textbf{número de interações e não tempo de especialista}.
\normalsize


O comportamento mais informativo dessa avaliação não é o valor médio, e sim a ocorrência de regressão. Em várias execuções o modelo piora ao longo da sequência de cliques: em um caso representativo, doze interações levaram a qualidade de 0,724 para 0,702. O especialista anota repetidamente e o modelo não avança, porque cada anotação nova reestima o banco de filtros e desfaz parte do que havia sido construído.

A proteção do banco não reduziu o número de cliques de forma significativa, e o confirmatório em imagens inéditas falseou tanto a hipótese quanto o mecanismo proposto para ela. A taxa de sucesso é reportada ao lado da contagem de cliques porque imagens que não atingem o alvo entram com o teto, e um braço poderia exibir contagem baixa simplesmente por desistir antes.

\begin{table}[htb]
\centering
\caption{Confirmatório em imagens inéditas.}
\label{tab_noc_conf}
% Cliques até a IoU alvo, por plasticidade do banco de filtros
% GERADO por flim_al/tabelas.py em 2026-10-08T14:01:07+00:00
% NAO EDITE A MAO: a proxima geracao sobrescreve, e o numero editado deixa de bater com evidencia/execucoes/.
% Proveniencia: noc_confirmatorio.proveniencia.json
\begin{adjustbox}{max width=\linewidth}
\begin{tabular}{lrrrrrrrrr}
\toprule
dataset & $\alpha$ & n & NoC médio & NoC mediano & taxa de sucesso & IoU final & $\Delta$ NoC vs $\alpha$=1 & p & McNemar \\
\midrule
Parasitas & 0.00 & 18 & 11.50 & 12.0 & 6\% & 0.5179 & 0.00 & 1.0000 & 0/0 (p=1.000) \\
Parasitas & 1.00 & 18 & 11.50 & 12.0 & 6\% & 0.5646 & --- & --- & --- \\
BraTS & 0.00 & 20 & 9.50 & 12.0 & 30\% & 0.2515 & -0.10 & 0.9138 & 2/2 (p=1.000) \\
BraTS & 1.00 & 20 & 9.60 & 12.0 & 30\% & 0.2770 & --- & --- & --- \\
\bottomrule
\end{tabular}
\end{adjustbox}

\end{table}

% Nota de noc_confirmatorio. GERADA por flim_al/tabelas.py.
\footnotesize NoC@X é o número de cliques até a IoU passar do alvo, com teto de 12; é o eixo de RITM (arXiv:2102.06583) e SimpleClick (arXiv:2210.11006), e \textbf{não} se compara com o NoC publicado deles, que usa outros conjuntos e outra definição de alvo. Imagem que não atinge o alvo entra com o teto, e por isso a \textbf{taxa de sucesso} aparece em toda linha: um braço pode mostrar NoC baixo por desistir mais cedo. \texttt{$\alpha$=1} é o FLIM puro, que refaz o banco de filtros a cada clique; \texttt{$\alpha$=0} congela o banco, e o clique chega ao modelo só pelos rótulos que o decoder \texttt{labeled\_marker} lê. Os testes são \textbf{separados por dataset}: agregar Schisto e BraTS, com 7\% e 30\% de sucesso, inflou um achado exploratório a p=0,0265 que separado não passa de p=0,10. O clique é simulado do ground truth pelo protocolo de Xu et al. (2016), com usuário que acerta sempre, então a tabela mede \textbf{número de interações e não tempo de especialista}.
\normalsize


Duas advertências de leitura acompanham a Tabela~\ref{tab_noc}. A primeira é que as linhas do conjunto de conjuntivite têm $n$ igual a um, e uma única execução não autoriza leitura alguma; elas constam para que a ausência de dado seja visível, e não para serem interpretadas. A segunda é de método, e vale mais que o resultado em si: uma versão anterior desta análise agregava os conjuntos de parasitas e BraTS numa única comparação e obtinha valor p de 0,0265, o que parecia um achado. Separados por conjunto, como manda o protocolo, nenhum dos dois passa de 0,10. A agregação inflava o resultado porque os conjuntos têm taxas de sucesso muito diferentes, de 7\% e 30\%, e misturá-las cria uma diferença que não é do critério.

\section{A intervenção no gerador de candidatos}

As seções anteriores convergem para um diagnóstico. Existe margem considerável na escolha da região. Os critérios de incerteza localizam o objeto mas pousam no interior. Os critérios de diversidade não localizam o objeto. E a anotação que os marcadores dos especialistas mostram ser valiosa, a que atravessa a fronteira, corresponde a uma fração mínima do conjunto de candidatos. Como um critério apenas ordena a lista que recebe, nenhuma escolha de escore poderia resolver a limitação.

A intervenção que decorre desse diagnóstico atua sobre o gerador. Em lugar de segmentar a imagem em superpixels, os candidatos passam a ser discos centrados em pontos do contorno previsto pelo modelo, de modo que atravessam a fronteira por construção. O contorno usado é o previsto, não o do ground truth, e portanto o procedimento é aplicável.

O experimento fixa o critério de seleção em sorteio nos dois braços e mantém o orçamento de anotação idêntico, verificado em mediana de trezentos pixels em ambos. A única diferença entre os braços é a origem da lista de candidatos, o que isola o efeito do gerador.

\begin{table}[htb]
\centering
\caption{Troca do gerador de candidatos, com critério e orçamento fixos.}
\label{tab_gerador}
% Gerador de candidatos: contorno previsto contra superpixel
% GERADO por flim_al/tabelas.py em 2026-10-08T14:01:07+00:00
% NAO EDITE A MAO: a proxima geracao sobrescreve, e o numero editado deixa de bater com evidencia/execucoes/.
% Proveniencia: gerador_contorno.proveniencia.json
\begin{adjustbox}{max width=\linewidth}
\begin{tabular}{lrrrrrrrrrrr}
\toprule
dataset & n & $\Delta$ médio & IC95\% & vitórias & p (t) & p (sinais) & p (perm.) & $\Delta$ pior caso & p  & $\Delta$ dispersão & p   \\
\midrule
Parasitas & 10 & 0.0772 & [0.0393, 0.1151] & 10 de 10 & 0.0013 & 0.0020 & 0.0020 & 0.2153 & 0.0001 & -0.1091 & 0.0000 \\
BraTS & 3 & -0.0089 & [-0.3263, 0.3084] & 1 de 3 & 0.9145 & 1.0000 & 0.7500 & -0.1063 & 0.2703 & 0.0464 & 0.1376 \\
Conjuntivite & 5 & 0.0452 & [-0.0452, 0.1357] & 4 de 5 & 0.2374 & 0.3750 & 0.2500 & 0.1759 & 0.0480 & -0.0577 & 0.1492 \\
\bottomrule
\end{tabular}
\end{adjustbox}

\end{table}

% Nota de gerador_contorno. GERADA por flim_al/tabelas.py.
\footnotesize $\Delta$ é a variação de F$\beta$ na validação ao acrescentar \textbf{uma} anotação, do braço de contorno menos o braço de superpixel. Os dois sorteiam o candidato e gastam os mesmos ~300 px: a \textbf{única} diferença é o gerador da lista, e por isso a comparação isola o gerador e não diz nada sobre qual escore usar. A unidade é a \textbf{semente}; os dois decodificadores de uma semente compartilham o encoder e entram colapsados por média. Pior caso e dispersão foram declarados no pré-registro, com o mecanismo escrito antes de rodar: o candidato de contorno é pequeno e local, acrescenta menos filtros, e o ganho marginal mede correlação negativa entre filtros acrescentados e ganho. O nulo do BraTS é \textbf{previsto}: a arquitetura dele fixa o banco em 8 filtros por camada, os dois braços acrescentam exatamente 24, e o deslocamento que a intervenção evita não pode ocorrer. Semente cuja base degenerou (F$\beta$=0) sai, porque $\Delta$ a partir de zero não pode ser negativo. \textbf{O desfecho primário vem com três testes}: t pareado, teste de sinais e permutação exata de sinais. Com 3 e 5 sementes a normalidade das diferenças não é verificável, e o t sozinho não autorizaria conclusão; onde os três concordam, como no conjunto de parasitas, a conclusão não depende da suposição. \textbf{ATENÇÃO ao pior caso da conjuntivite}: o valor de 0,0480 é bruto e \textbf{não sobrevive} à análise de sensibilidade que casa o número de candidatos entre os braços, na qual passa a 0,1032. O gerador de contorno entregou UM único candidato naquela semente, e mínimo de um sorteio contra mínimo de oito não é comparação. Esse desfecho \textbf{não deve ser lido como significativo}.
\normalsize


No conjunto de parasitas o efeito é claro: ganho médio de 0,0772 de $F_\beta$, com intervalo de confiança entre 0,0393 e 0,1151, valor p de 0,0013, e vantagem em todas as dez repetições. O resultado sobrevive à correção de Benjamini-Hochberg aplicada sobre os três conjuntos de dados. Os dois desfechos secundários, declarados no pré-registro com o mecanismo escrito antes da execução, acompanham: o pior caso melhora 0,2153 e a dispersão entre candidatos cai 0,1091.

O comportamento nos outros dois conjuntos merece atenção. Na conjuntivite o efeito vai na mesma direção, com vantagem em quatro das cinco repetições, mas sem significância no desfecho primário. No BraTS o efeito desaparece, e esse resultado nulo é previsto pelo mecanismo: a arquitetura empregada nesse conjunto fixa o banco em oito filtros por camada, de modo que ambos os braços acrescentam exatamente o mesmo número de filtros e o deslocamento que a intervenção evita não pode ocorrer. O conjunto funciona, assim, como controle do mecanismo.

Duas limitações do gerador foram medidas e devem acompanhar qualquer leitura do resultado. A primeira é que ele exige uma predição não trivial: quando a máscara prevista é vazia ou cobre a imagem inteira, não há contorno e o conjunto de candidatos fica vazio, enquanto o gerador por superpixels funciona em qualquer situação. A segunda é que os candidatos gerados apresentam fração de objeto próxima de um, isto é, caem no interior do objeto verdadeiro e não sobre sua fronteira, porque o modelo subsegmenta e o contorno previsto se situa dentro do contorno real. O ganho ocorre apesar disso, o que sugere que parte do efeito decorre de o candidato ser pequeno e localizado, acrescentando menos filtros ao banco, e não apenas de sua posição. Separar essas duas contribuições exige um experimento adicional, discutido no Capítulo \ref{cap_5}.

\subsection{Análise de sensibilidade dos desfechos secundários}
\label{sec_sensibilidade}

Uma terceira limitação apareceu na conferência dos dados e afeta a leitura dos secundários, não do desfecho primário. O desenho declarou oito candidatos por braço, mas o gerador de contorno entregou menos em três sementes: três no conjunto de parasitas, três no BraTS e \emph{uma} na conjuntivite. O gerador por superpixels entregou oito em todas.

A consequência não é a mesma para os três desfechos. A média é estimador não viesado da média do braço qualquer que seja o número de candidatos, de modo que o desfecho primário não se move. Mas o desfecho de pior caso é um \emph{mínimo}, e o mínimo de oito sorteios é, por construção, mais extremo que o mínimo de um, mesmo que a distribuição subjacente seja idêntica. O viés empurra na direção que a tabela reporta, o que obriga a medi-lo.

Refazendo o teste com o braço de superpixels subamostrado ao mesmo número de candidatos, sem reposição, com duas mil repetições e o mesmo teste pareado por semente, obtém-se o seguinte. No conjunto de parasitas o pior caso passa de 0,2153 para 0,2085, com valor p de 0,0002, e a dispersão passa de $-0{,}1091$ para $-0{,}1078$, com valor p abaixo de 0,0001: ambos sobrevivem, e o efeito não era artefato do número de candidatos. Na conjuntivite, porém, o pior caso passa de 0,1759 com valor p de 0,0480 para 0,1369 com valor p de 0,1032, isto é, \textbf{cruza o limiar e deixa de ser significativo}. A semente responsável entregou um único candidato de contorno, e mínimo de um sorteio contra mínimo de oito não é comparação.

O desfecho secundário de pior caso na conjuntivite não é, portanto, apresentado como significativo nesta dissertação. Esta análise é de sensibilidade sobre um desfecho pré-registrado, e não um desfecho novo.

\subsection{O critério sobre a lista corrigida}

A intervenção acima troca o gerador mantendo o escore fixo em sorteio, de propósito, para isolar uma coisa de cada vez. Resta a pergunta complementar: uma vez corrigida a lista, o escore volta a ter utilidade?

Uma campanha adicional no conjunto de parasitas, com dez sementes completas, indica que sim, e os dois efeitos se somam de forma quase exata. Aplicar o critério de entropia sobre candidatos de contorno rende 0,0262 de $F_\beta$ sobre sortear dentro da mesma lista, com vantagem em nove das dez sementes e valor p de 0,0185. A troca do gerador isolada rende os 0,0772 já relatados. E o efeito combinado, critério sobre contorno contra sorteio sobre superpixels, rende 0,1033, com vantagem nas dez sementes e valor p de 0,0002. A soma das partes, 0,1034, coincide com o efeito total medido.

A aditividade é coerente com o diagnóstico: o gerador resolve \emph{o que está na lista} e o escore resolve \emph{qual item da lista}, e são operações sobre dimensões distintas do problema.

Duas ressalvas limitam o alcance desse resultado. Ele vem de um único conjunto de dados, e \textbf{não houve pré-registro para estas comparações}, de modo que elas são exploratórias e servem para orientar o próximo passo, não para afirmar. Registra-se também que a terceira linha dessa campanha reproduz exatamente o ganho de 0,0772 já publicado, com o mesmo valor p e as mesmas dez vitórias, o que serve de verificação interna: um cálculo novo, sobre uma campanha nova, devolveu o número anterior.

\section{Comparação com redes treinadas por retropropagação}

Para situar a ordem de grandeza dos valores acima, cinco redes de segmentação treinadas por retropropagação foram avaliadas nos dois conjuntos de maior porte, com três partições cada. No conjunto de parasitas os valores de $F_\beta$ ficam entre 0,3763 e 0,4450, e no BraTS entre 0,7421 e 0,8511.

A comparação direta com os números do \ac{FLIM} apresentados neste capítulo \textbf{não é válida} e não é feita aqui, por três razões que convém enunciar em vez de contornar. As redes por retropropagação foram treinadas com o conjunto de treino completo e anotação densa, enquanto o \ac{FLIM} opera com três a cinco imagens e traços esparsos, de modo que o orçamento de anotação difere por ordens de grandeza. Os protocolos de avaliação e as partições não coincidem. E, principalmente, esses resultados \textbf{não estão no registro canônico de execuções} que sustenta todas as demais tabelas deste capítulo, tendo proveniência mais fraca que o restante.

O que eles autorizam é apenas uma observação de contexto: no conjunto de parasitas, o patamar alcançado por redes com supervisão densa é da ordem de 0,44, enquanto o \ac{FLIM} com três marcadores reais atinge 0,744 no decodificador \texttt{FLIM\_lm} sob o protocolo de avaliação adotado aqui. Os dois números medem coisas diferentes e a diferença de protocolo é suficiente para explicar a distância; o que não se sustenta é a suposição de que supervisão densa seja condição necessária para desempenho utilizável neste domínio.

\section{Resultados de natureza metodológica}

Três resultados não dizem respeito ao desempenho dos métodos, mas à validade das medições, e por isso são relatados separadamente.

\subsection{Reprodutibilidade do treinamento}

O treinamento do codificador não era reprodutível entre execuções. Duas chamadas idênticas, com os mesmos arquivos de marcadores e no mesmo processo, produziam pesos diferentes quando o número de candidatos a filtro ultrapassava o limite da arquitetura: a diferença máxima observada nos pesos da segunda camada foi de 0,134. Abaixo desse limite o procedimento devolve os próprios candidatos sem agrupá-los, e por isso é estável; acima dele o agrupamento precisa escolher entre muitos candidatos, e a ordem de redução das somas em bibliotecas de álgebra linear com múltiplas linhas de execução altera o resultado no último dígito, o que basta para o agrupamento convergir para outra configuração. O efeito se amplifica entre camadas, já que cada camada recebe como entrada as ativações da anterior.

Fixando as bibliotecas em uma única linha de execução, a diferença passa a ser exatamente zero nos dois regimes. Todas as campanhas posteriores a essa descoberta foram executadas sob essa restrição.

\subsection{Defeitos que invalidaram medições anteriores}

Três defeitos de implementação foram encontrados e corrigidos, e os números produzidos antes das correções não são válidos. O procedimento de extração de atributos para os critérios de diversidade assumia imagens de três bandas, o que impedia sua execução no conjunto BraTS, de banda única. A leitura de imagens assumia extensão fixa, o que impedia a execução no conjunto de conjuntivite. E o filtro de área do pós-processamento estava configurado com a faixa adequada ao conjunto de parasitas, cujos objetos são pequenos, o que eliminava toda componente prevista nas imagens de conjuntivite: nenhum dos objetos examinados cabia na faixa, e o $F_\beta$ mediano passou de 0,024 para 0,594 após a correção.

Esse último caso tem interesse além do defeito pontual. Ele indica que os parâmetros de pós-processamento publicados não transferem entre domínios, e que aplicá-los sem reavaliação pode produzir resultados sem relação com o método avaliado.

\subsection{Falsos positivos capturados pelo protocolo}

O protocolo de pré-registro, descrito no Capítulo \ref{cap_3}, capturou quatro resultados que não se sustentaram. Dois deles haviam atingido significância em análises exploratórias e foram submetidos a confirmatórios com repetições inéditas: um apresentava valor p de 0,0016 e passou a 0,586, outro apresentava 0,082 e passou a 0,567. Um terceiro desapareceu quando o número de repetições foi completado, com o valor p passando de 0,0265 para 0,1385. O quarto foi falseado não apenas na hipótese como no mecanismo proposto, pelo critério declarado antes da execução.

\begin{table}[htb]
\centering
\caption{Composição do registro de execuções.}
\label{tab_proveniencia}
% Proveniência disponível por família de experimento
% GERADO por flim_al/tabelas.py em 2026-10-08T14:01:07+00:00
% NAO EDITE A MAO: a proxima geracao sobrescreve, e o numero editado deixa de bater com evidencia/execucoes/.
% Proveniencia: proveniencia_registro.proveniencia.json
\begin{tabular}{lrrr}
\toprule
família & execuções & sem seed & sem commit \\
\midrule
ganho\_marginal & 2448 & 0\% & 0\% \\
al\_corrected\_results & 2316 & 42\% & 100\% \\
tabela\_k\_por\_modelo & 1840 & 0\% & 0\% \\
final\_comparison\_v3 & 1029 & 4\% & 100\% \\
benchmark\_conjunctiva & 588 & 68\% & 100\% \\
sel\_cs\_grande & 540 & 0\% & 100\% \\
al\_backprop\_results & 492 & 20\% & 100\% \\
paper\_selection\_pac2 & 432 & 0\% & 100\% \\
paper\_selection\_POOL\_COM\_VAZAMENTO & 360 & 0\% & 100\% \\
paper\_selection\_v2 & 324 & 0\% & 100\% \\
paper\_selection\_v3 & 324 & 0\% & 100\% \\
paper\_selection\_rand & 321 & 0\% & 100\% \\
densidade & 262 & 0\% & 0\% \\
cl\_plasticidade & 254 & 0\% & 0\% \\
sel\_degen & 216 & 0\% & 100\% \\
sel\_medoide & 216 & 0\% & 100\% \\
sel\_proposto & 216 & 0\% & 100\% \\
final\_comparison\_v2 & 210 & 10\% & 100\% \\
il\_noc & 207 & 0\% & 0\% \\
clique & 183 & 0\% & 0\% \\
paper\_selection\_coreset & 162 & 0\% & 100\% \\
onde\_marcar & 150 & 0\% & 0\% \\
final\_comparison & 147 & 57\% & 100\% \\
al\_encoder\_results\_acq\_comparison & 144 & 44\% & 100\% \\
orcamento & 144 & 100\% & 100\% \\
artigo\_vs\_regiao & 130 & 0\% & 0\% \\
al\_curve & 84 & 100\% & 100\% \\
brats & 84 & 25\% & 100\% \\
al\_encoder\_results & 70 & 70\% & 100\% \\
al\_encoder\_results\_dt & 62 & 47\% & 100\% \\
paper\_selection\_brats & 54 & 0\% & 100\% \\
paper\_selection\_uB\_POOL\_COM\_VAZAMENTO & 54 & 0\% & 100\% \\
al\_bald\_results & 48 & 44\% & 100\% \\
al\_flim\_curve & 42 & 100\% & 100\% \\
curva & 24 & 100\% & 100\% \\
al\_region\_results & 16 & 44\% & 100\% \\
\bottomrule
\end{tabular}

\end{table}

% Nota de proveniencia_registro. GERADA por flim_al/tabelas.py.
\footnotesize Diagnóstico, não resultado. Linha com 100\% sem seed não sustenta afirmação sobre variância entre execuções.
\normalsize


\section{Síntese}

Os resultados sustentam a seguinte leitura, em quatro elos, cada um medido.

\textbf{Primeiro, a seleção de imagem não é a alavanca.} Nenhum critério de \ac{AL} apresenta superioridade sobre a escolha aleatória que sobreviva à correção para múltiplas comparações, e o procedimento supervisionado do trabalho de referência, que lê o gabarito de todo o conjunto, também não se destaca. Isso não é demonstração de equivalência: é ausência de evidência de superioridade nas condições avaliadas.

\textbf{Segundo, a posição da anotação é.} Com orçamento de pixels fixo e nas mesmas imagens, a melhor região examinada supera a sorteada por 0,1351 a 0,3089 de $F_\beta$, nas seis células, e a pior fica 0,2075 a 0,4954 abaixo. O prêmio e o risco são da mesma ordem, o que significa que anotar mais não é, por si, melhorar.

\textbf{Terceiro, os critérios capturam parte dessa margem, e a parte que escapa é o custo do método.} A lacuna entre o teto e o que o critério entrega vai de 0,0839 a 0,2667 e é significativa em quatro das seis células. O desempenho do escore depende do decodificador e do domínio: há uma célula em que ele captura 84\% da margem e outra em que captura praticamente nada.

\textbf{Quarto, o gargalo está na lista e não no escore.} Regiões que atravessam a fronteira do objeto, que são as que a comparação com os marcadores reais mostra terem valor, representam de 2\% a 12\% dos candidatos produzidos por superpixels. Como um critério apenas ordena a lista que recebe, nenhuma escolha de escore resolveria isso. A intervenção derivada desse diagnóstico, que troca o gerador mantendo escore e orçamento fixos, produz ganho de 0,0772 de $F_\beta$ no conjunto de parasitas, em dez de dez repetições, sobrevivendo à correção; o resultado nulo no BraTS foi previsto pelo mecanismo antes da execução e funciona como controle dele.

A razão estrutural que percorre os quatro elos é a mesma: no \ac{FLIM}, a anotação adicional \emph{reestima} o extrator de atributos em vez de refiná-lo, e o banco de filtros é compartilhado por todas as imagens, de modo que o acréscimo é uma intervenção global e de risco. O Capítulo \ref{cap_5} discute o que esses resultados autorizam afirmar e o que permanece em aberto.
